\chapter{附录：补充公式与统计说明}

\section{固定组采样下的 signal collapse 概率}

设单个 prompt 的单次答对概率为 $p_x$，组大小为 $n$。若组内全部采样均正确，则其概率为 $p_x^n$；若组内全部采样均错误，则概率为 $(1-p_x)^n$。因此，发生 signal collapse 的 prompt 级概率为
\begin{equation}
    P_{\mathrm{collapse}}(x;n)=p_x^n+(1-p_x)^n.
\end{equation}

若对训练分布求期望，即可得到总体坍塌率
\begin{equation}
    P_{\mathrm{collapse}}(n)=\mathbb{E}_{x\sim d_0}[p_x^n+(1-p_x)^n].
\end{equation}

该表达式说明，只有当 prompt 级正确率大量集中在中间区间且组大小足够大时，坍塌率才会明显下降；一旦训练集中同时存在大量接近 $0$ 或接近 $1$ 的样本，固定小组采样就会持续暴露全错或全对问题。

\section{单位预算信息效率近似}

对二值奖励而言，单样本方差信息与 $p_x(1-p_x)$ 成正比。若进一步考虑训练后期获取负样本的期望成本，可得到单位预算信息效率近似
\begin{equation}
    \eta(p_x)\propto p_x(1-p_x)^2.
\end{equation}

当 $p_x=0.9$ 时，$\eta(p_x)=0.009$；当 $p_x=0.3$ 时，$\eta(p_x)=0.147$。这说明中等偏难样本在相同预算下更可能产出高价值训练信号，也为本文采用难度感知预算分配提供了直观解释。
